\section{Observability statement and scope}\label{s0.-observability-statement-and-scope}

The main article concerns the exact matched-parameter signature \(\Sigma=(G,U_{xx},\operatorname{Tr}S)\), where \(G\) is the centered spatial growth class. For \(M>0\), \(\rho>0\), and \(\mathrm{Pe}\ge0\), it is injective on the seven nonempty fixed-zero reset protocols. Its proof has three branches: spatial growth recovers position reset, \(U_{xx}\) recovers orientation reset, and positive kinetic-trace gaps recover velocity reset. It does not assert global minimality. Sections S1--S7 retain the full stationary/asymptotic second-order record as the supplementary completeness result. Its heterogeneous stationary and asymptotic entries are not asserted to be one finite-time record.

This supplement gives the generator, closed second-order hierarchy, exhaustive pairwise proof, catalogue of restricted observational identities, and physical limiting regimes for the two-dimensional inertial active Brownian particle with deterministic Poisson resetting. It is written to be read independently.

\subsection{Numerical corroboration and trajectory-study disclosure}\label{s0.1-numerical-corroboration}

The moment results are exact consequences of the closed hierarchy, and the internal-sector invariance follows from the reset generator; both were obtained symbolically. An independently implemented finite-grid moment calculation, coded from the model rather than from the closed-form solutions, reproduced the exact formulas at its prespecified evaluation points; it is corroborative and does not prove continuum-wide inequalities.

The immutable historical S-071 and S-073 stochastic trajectory datasets were analyzed separately under their prespecified sampling plans, without pooling or further sampling. S-071 validated 14 of 46 items, left 32 unresolved, and had 0 contradicted; S-073 validated 32 of 46, left 14 unresolved, and had 0 contradicted. All 14 S-073 unresolved rows missed at least one precision requirement, and seven also failed containment. Their original criteria were $q\,\mathrm{SE}\leq0.5\epsilon$ and $q\,\mathrm{SE}+B\leq0.75\epsilon$. Neither study supports a quantitative trajectory-validation claim; complete records remain in the reproducibility archive.

\subsection{Independent trajectory simulation}\label{s0.2-s074-trajectory-method}

Study S-074 used 23,519,232 independent production trajectories in 41 parameter--protocol cells, with 49,152--1,146,880 trajectories per cell. Table~\ref{tab:s074} retains all 80 planned comparisons. After inspection of the raw-finest results, Richardson extrapolation was adopted using the existing coupled trajectories, without new sampling, pooling, row selection or top-up. The raw-finest analysis contained 80/80 intervals, with 0 contradicted and 78/80 meeting the stricter target; its $\operatorname{Tr}S$ residuals reached about $+8.1$ SE and kinetic-trace convergence ratios were 1.999--2.006. Richardson did not rescue containment. The original estimates remain archived; the seven raw kinetic-trace estimates lie 0.10--0.27\% above the exact references. All 80 paired SEs are below 1\% of their declared precision scales (maximum 0.85858\%), and all extrapolated intervals are contained.

Trajectories start from zero with full translational and rotational noise. Exact exponential reset times are shared across levels; accumulated coarse Brownian increments are flushed before each reset and observation. For an uninterrupted Euler--Maruyama step of length $\delta$, put $a_\delta=1-\delta/M$ and $b_\delta=\mathrm{Pe}\delta/M$. The velocity moment update is
\[
E|v'|^2=a_\delta^2E|v|^2+2a_\delta b_\delta E[v\cdot u]+b_\delta^2+4\delta/M^2.
\]
The angular Brownian update and reset action are exact at segment endpoints. These smooth between-event moment maps support the model-specific first-order weak expansion $E[X_h]=X_\star+c_1h+O(h^2)$ for a fixed observation-window statistic; the paired combination $R=2X_h-X_{2h}$ then cancels the leading bias. Centering and fixed-window slope fitting are performed at each level before combination. This is not a proof of a uniform stationary expansion or a rigorous global error bound. For $U_{xx}$, the coupled angular values agree at recording times up to roundoff, so the leading bias may vanish; a convergence ratio near two is not required for those rows.

The finest archived step is $h=1/512$, $1/1024$, or $1/2048$ according to the archive rule, with summed Brownian increments pairing $h$, $2h$ and $4h$. For each row and late-time window, the analysis also evaluates $R'=2X_{2h}-X_{4h}$ using stored joint covariance. The reported discrepancy interval is $R_0-m_{\mathrm{ref}}\ \mathbin{\pm}\ [q\,\mathrm{SE}_R+B_R+B_{\mathrm{late}}]$, where $B_R=|2X_h-3X_{2h}+X_{4h}|+q\,\mathrm{SE}(2X_h-3X_{2h}+X_{4h})$ and $B_{\mathrm{late}}=\max(\text{original primary envelope},\max_w|R_w-R_0|,\text{original transient residual})$. Late-time error is assessed separately from step bias and sampling SE. Each cell used 65 equally spaced records to $T=64$ or $160$, with primary window $[T/2,T]$ and sensitivity windows $[2T/3,T]$ and $[3T/4,T]$.

An excluded pilot of 4096 trajectories per cell determined fixed fresh sample sizes from SEs alone, inflated by 1.20 and rounded up to batches of 8192. Neither pilot nor historical data enter production estimates. Linear moments use one late-window summary per independent trajectory. Centered covariances and diffusion slopes use a whole-trajectory delete-one jackknife with recentering and slope refitting; pilot jackknife/bootstrap SE ratios were 0.95465--1.05064. Correlated observation times are not treated as independent samples.

The sampling-precision scale $C_{\mathrm{precision}}$ is $|m_{\mathrm{ref}}|$ for a nonzero reference, including small diffusion coefficients, and $\sqrt{S_{xx,\mathrm{ref}}}$ for a symmetry-zero mean velocity. The agreement margin is separately defined by $\epsilon=0.005C_{\mathrm{char}}+0.01|m_{\mathrm{ref}}|$, with no relative term for symmetry zeros. Here $C_{\mathrm{char}}=|m_{\mathrm{ref}}|$ for $U_{xx}$, $\operatorname{Tr}S$, $\operatorname{Tr}R$ and stationary centered spatial variance; $\max(1,|D_{\mathrm{eff,ref}}|)$ for diffusion; $\sqrt{S_{xx,\mathrm{ref}}}$ for mean velocity; and $\sqrt{\operatorname{Tr}R_{\mathrm{ref}}\operatorname{Tr}S_{\mathrm{ref}}}$ for $r\cdot v$. For $V\Theta$ at $M=2$, $\mathrm{Pe}=3$, $\rho=6$, the SE is 0.85858\% of the reference and the actual discrepancy is 1.24\%, whereas $\epsilon=0.0051507476$ permits about 34.2\% of the reference. Sampling precision and agreement tolerance are not interchangeable.

Both raw and extrapolated analyses miss the stricter original target $\tau=\min(0.01C_{\mathrm{precision}},\epsilon/(3q))$ for two stationary-variance rows at $M=2$, $\mathrm{Pe}=3$, $\rho=0.5$: the Richardson SE exceeds $\tau$ as $0.0238650088>0.0226464282$ for $PV$ and $0.0315097096>0.0310579587$ for $P\Theta$. Thus 78 of 80 rows meet this stricter target. At the core row, the internally identical $P\Theta$ kinetic trace is $+3.01$ SE while $\Theta$ is about $-0.5$ SE; this is compatible with sampling fluctuation, but does not establish its cause. Table~S1's V denotes containment of the expanded discrepancy interval within the declared tolerance band, not passage of every precision requirement or universal relative accuracy. The multiplier $q=2.5758293035$ gives nominal pointwise 99\% sampling intervals; the systematic allowances are empirical assessments, not rigorous bounds or simultaneous 99\% coverage.

\input{PHASEMAP_S074_Table.tex}

\section{Model, reset semantics, and generator}\label{s1.-model-reset-semantics-and-generator}

Let \(r=(x,y)^\mathsf T\), \(v=(v_x,v_y)^\mathsf T\), and \(u(\theta)=(\cos\theta,\sin\theta)^\mathsf T\). Between resets,

\[
dr=v\,dt,\qquad
M\,dv=-(v-\mathrm{Pe}\,u)\,dt+\sqrt{2}\,dW_t,\qquad
d\theta=\sqrt{2}\,dW_{r,t},
\]

where \(W_t\) is a standard two-component Wiener process, \(W_{r,t}\) is an independent scalar Wiener process, and

\[
M>0,\qquad \mathrm{Pe}\ge 0,\qquad \rho\ge 0.
\]

Reset times are those of an independent Poisson process of rate \(\rho\). Each reset is instantaneous. With \(0\) denoting the appropriate scalar or two-component zero, the seven nonempty deterministic reset protocols are

\[
\begin{array}{c@{\quad}c@{\qquad}c@{\quad}c}
P&(0,v,\theta)&V&(r,0,\theta)\\
\Theta&(r,v,0)&PV&(0,0,\theta)\\
P\Theta&(0,v,0)&V\Theta&(r,0,0)\\
PV\Theta&(0,0,0)&&
\end{array}
\]

Every unlisted state variable is retained exactly. For a smooth observable \(f(r,v,\theta)\), the backward generator is

\[
\boxed{
\begin{aligned}
\mathcal L_q f
={}&v\mathbin{\cdot}\nabla_r f
-\frac{v-\mathrm{Pe}\,u}{M}\mathbin{\cdot}\nabla_v f
+\frac{1}{M^2}\Delta_v f+\partial_\theta^2 f\\
&+\rho\left[f(R_q(r,v,\theta))-f(r,v,\theta)\right].
\end{aligned}}
\]

The common default initial state is \(r(0)=0\), \(v(0)=0\), and \(\theta(0)=0\). Thus \(\bar r(0)=\bar v(0)=0\), \(\bar u(0)=e=(1,0)^\mathsf T\), all raw second moments containing \(r\) or \(v\) initially vanish, and \(U(0)=ee^\mathsf T\).



\section{Temporal scope of the second-order record}\label{s2.-temporal-scope-of-the-second-order-record}

For a protocol \(q\), define

\[
\mathcal R_2(q;M,\mathrm{Pe},\rho)
\]

to be the fixed stationary-or-asymptotic component-level output used by the classification. For a position-reset protocol, the record contains the stationary long-time limits of all means and raw second moments below, the centered spatial covariance, and the finite stationary variance. For a protocol without position reset, the record contains the stationary internal \((v,u)\) moments and the long-time spatial polynomial coefficients: \(\bar r(t)=\bar v\,t+b+o(1)\), \(\operatorname{Cov}r(t)=2Dt+O(1)\), the corresponding raw \(R,C,Q\) coefficients, \(D_{\mathrm{eff}}\), and the centered and raw spatial-growth classes. Equality means equality of every applicable stationary value and every applicable long-time coefficient at matched \((M,\mathrm{Pe},\rho)\). The theorem does not assert equality or distinguishability of unspecified finite-time transients, and it does not claim finite-sample recoverability.

The raw matrices used throughout are

\[
\begin{gathered}
R=E[rr^\mathsf T],\quad S=E[vv^\mathsf T],\quad
C=C_{rv}=E[rv^\mathsf T],\\
Q=E[ru^\mathsf T],\quad W=E[vu^\mathsf T],\quad
U=E[uu^\mathsf T],
\end{gathered}
\]

with means \(\bar r=E[r]\), \(\bar v=E[v]\), and \(\bar u=E[u]\). The required named scalar contractions are

\[
\mathrm{TrR}:=\operatorname{Tr}R=E|r|^2,\qquad
\mathrm{TrS}:=\operatorname{Tr}S=E|v|^2,\qquad
\mathrm{TrC}:=\operatorname{Tr}C=E[r\mathbin{\cdot}v].
\]

The centered spatial covariance is \(\operatorname{Cov}(r)=R-\bar r\bar r^\mathsf T\). For a diffusive protocol,

\[
D_{\mathrm{eff}}
=\lim_{t\to\infty}\frac{\operatorname{Tr}\operatorname{Cov}(r(t))}{4t}.
\]

\section{Ordered redundant 28-coordinate frame and exact linear moment hierarchy}\label{s3.-ordered-redundant-28-coordinate-frame-and-exact-linear-moment-hierarchy}

We use the ordered coordinate frame

\[
\begin{aligned}
\mathcal F_2=\big(&1;\\
&r_x,r_y;\\
&v_x,v_y;\\
&u_x,u_y;\\
&R_{xx},R_{xy},R_{yy};\\
&S_{xx},S_{xy},S_{yy};\\
&C_{xx},C_{xy},C_{yx},C_{yy};\\
&Q_{xx},Q_{xy},Q_{yx},Q_{yy};\\
&W_{xx},W_{xy},W_{yx},W_{yy};\\
&U_{xx},U_{xy},U_{yy}\big).
\end{aligned}
\]

These are 28 stored coordinates, not an unrestricted 28-element basis. Every physical moment record lies on the invariant manifold

\[
\boxed{\texttt{one}=1,\qquad
\operatorname{Tr}U=U_{xx}+U_{yy}=1,}
\]

or, in homogeneous coordinate form, \(\texttt{one}-U_{xx}-U_{yy}=0\). The symmetric matrices \(R,S,U\) retain three stored components, while the nonsymmetric matrices \(C,Q,W\) retain all four. Together with the constant and six first moments, the frame contains \(1+6+9+12=28\) entries. The exact trace constraint makes this a redundant moment frame; retaining the redundancy keeps the component bookkeeping explicit.

In component notation, the same exact order is

\begin{verbatim}
one,
r_x, r_y,
v_x, v_y,
u_x, u_y,
rr_xx, rr_xy, rr_yy,
vv_xx, vv_xy, vv_yy,
rv_xx, rv_xy, rv_yx, rv_yy,
ru_xx, ru_xy, ru_yx, ru_yy,
vu_xx, vu_xy, vu_yx, vu_yy,
uu_xx, uu_xy, uu_yy
\end{verbatim}

Here \texttt{rr}, \texttt{vv}, \texttt{rv}, \texttt{ru}, \texttt{vu}, and \texttt{uu} denote \(R,S,C,Q,W,U\), respectively.

Write the reset indicators of protocol \(q\) as \((p,\nu,\eta)\in\{0,1\}^3\), for position, velocity, and orientation reset, respectively. For binary indicators, write \(a\vee b=a+b-ab\). Direct application of \(\mathcal L_q\) gives the exact closed hierarchy

\[
\dot 1=0,
\]

\[
\begin{aligned}
\dot{\bar u}
&=-\bar u+\rho\eta(e-\bar u),\\
\dot U
&=2\operatorname{Tr}(U)I-4U+\rho\eta(ee^\mathsf T-U),
\end{aligned}
\]

This trace form is the exact coordinate representation used by the 28-component closure: \(\dot U_{xx}=-2U_{xx}+2U_{yy}+\rho\eta(1-U_{xx})\), \(\dot U_{xy}=-(4+\rho\eta)U_{xy}\), and \(\dot U_{yy}=2U_{xx}-2U_{yy}-\rho\eta U_{yy}\). The physical invariant \(\operatorname{Tr}U=1\) then also permits the equivalent shorthand \(2I-4U\) on admissible states. Indeed, taking the trace gives

\[
\frac{d}{dt}\big(\operatorname{Tr}U-1\big)
=-\rho\eta\big(\operatorname{Tr}U-1\big).
\]

Protocols without orientation reset (\(\eta=0\)) preserve the trace constraint, while protocols with orientation reset (\(\eta=1\)) restore any trace deviation exponentially at rate \(\rho\). Thus every one of the seven generators leaves the physical invariant manifold unchanged, and all physical conclusions below are restricted to that manifold.

\[
\begin{aligned}
\dot{\bar v}
&=-\frac{\bar v-\mathrm{Pe}\,\bar u}{M}-\rho\nu\bar v,\\
\dot W
&=\frac{\mathrm{Pe}}{M}U-\left(1+\frac1M\right)W
-\rho(\nu\vee\eta)W
+\rho\eta(1-\nu)\bar v e^\mathsf T,\\
\dot S
&=-\frac{2}{M}S+\frac{\mathrm{Pe}}{M}(W+W^\mathsf T)
+\frac{2}{M^2}I-\rho\nu S,
\end{aligned}
\]

and

\[
\begin{aligned}
\dot{\bar r}
&=\bar v-\rho p\bar r,\\
\dot Q
&=W-Q-\rho(p\vee\eta)Q
+\rho\eta(1-p)\bar r e^\mathsf T,\\
\dot C
&=S-\frac1M C+\frac{\mathrm{Pe}}{M}Q
-\rho(p\vee\nu)C,\\
\dot R
&=C+C^\mathsf T-\rho pR.
\end{aligned}
\]

These tensor equations are precisely the component equations \(\dot m_q=A_qm_q\) in the ordered frame \(\mathcal F_2\). The constant entry carries the affine terms. No moment outside \(\mathcal F_2\) is generated. Any uniqueness statement for this coordinate system is restricted to the physical invariant manifold above; no unrestricted ambient-space uniqueness is claimed.

\subsection{Long-time solution for position-reset protocols}\label{s3.1-long-time-solution-for-position-reset-protocols}

For \(P,PV,P\Theta,PV\Theta\) with \(\rho>0\), every position block has a strict decay contribution. The hierarchy is triangular from \((\bar u,U)\) through \((\bar v,W,S)\) to \((\bar r,Q,C,R)\), so it has a unique finite stationary solution on the physical invariant manifold.

The orientation block is

\[
\begin{array}{c|cc}
q & \bar u & U\\ \hline
P,PV & 0 & I/2\\[2mm]
P\Theta,PV\Theta
& \dfrac{\rho}{1+\rho}e
& \operatorname{diag}\!\left(\dfrac{\rho+2}{\rho+4},
                              \dfrac{2}{\rho+4}\right).
\end{array}
\]

The velocity blocks are

\[
\begin{array}{c|cc}
q & \bar v & W\\ \hline
P
& \mathrm{Pe}\,\bar u
& \dfrac{\mathrm{Pe}\,U}{1+M}\\[2mm]
PV
& \dfrac{\mathrm{Pe}\,\bar u}{1+M\rho}
& \dfrac{\mathrm{Pe}\,U}{1+M(1+\rho)}\\[2mm]
P\Theta
& \mathrm{Pe}\,\bar u
& \dfrac{\mathrm{Pe}\,U+M\rho\,\bar v e^\mathsf T}
        {1+M(1+\rho)}\\[3mm]
PV\Theta
& \dfrac{\mathrm{Pe}\,\bar u}{1+M\rho}
& \dfrac{\mathrm{Pe}\,U}{1+M(1+\rho)}.
\end{array}
\]

For all four protocols,

\[
\begin{aligned}
\bar r&=\frac{\bar v}{\rho},\\
S&=\frac{\mathrm{Pe}(W+W^\mathsf T)+2I/M}
        {2+M\rho\,\mathbf 1_{\{V\ {\rm reset}\}}},\\
Q&=\frac{W}{1+\rho},\\
C&=\frac{MS+\mathrm{Pe}\,Q}{1+M\rho},\\
R&=\frac{C+C^\mathsf T}{\rho}.
\end{aligned}
\]

Consequently, all four position-reset protocols are stationary/localized in centered position and have bounded raw MSD.



\subsection{Long-time solution for non-position protocols}\label{s3.2-long-time-solution-for-non-position-protocols}

The internal \((v,u)\) blocks of \(V,\Theta,V\Theta\) are exactly the same as those of \(PV,P\Theta,PV\Theta\), respectively. Define

\[
\Sigma_u=U-\bar u\bar u^\mathsf T,\qquad
\Sigma_v=S-\bar v\bar v^\mathsf T,\qquad
K=W-\bar v\bar u^\mathsf T.
\]

With

\[
\alpha=
\begin{cases}
1,&q=V,\\
1+\rho,&q=\Theta,V\Theta,
\end{cases}
\qquad
B=
\begin{cases}
1+M\rho,&q=V,V\Theta,\\
1,&q=\Theta,
\end{cases}
\]

the centered long-time transport law is

\[
\lim_{t\to\infty}\operatorname{Cov}(r,u)=\frac{K}{\alpha},
\qquad
\lim_{t\to\infty}\operatorname{Cov}(r,v)
=\frac{M\Sigma_v+\mathrm{Pe}\,K/\alpha}{B},
\]

\[
\operatorname{Cov}(r(t))=2Dt+O(1),\qquad
D=\frac{\operatorname{Cov}(r,v)+
        \operatorname{Cov}(r,v)^\mathsf T}{2},
\qquad
D_{\mathrm{eff}}=\frac{\operatorname{Tr}D}{2}.
\]

Thus all three non-position protocols are diffusive in centered covariance. Conditioned on the reset clock, the thermal and active velocity parts are independent and the thermal part has zero mean; hence their cross-covariance vanishes and \(D=D_{\mathrm{th}}+D_{\mathrm{act}}\), with \(D_{\mathrm{act}}\succeq0\). For velocity-reset protocols \(D_{\mathrm{th}}=2I/[(1+M\rho)(2+M\rho)]\), while for \(\Theta\), \(D_{\mathrm{th}}=I\). Thus \(D\) is positive definite throughout the domain. For \(V\), \(\bar v=0\), so raw MSD is also diffusive. For the two orientation-reset protocols,

\[
\bar v_\Theta=\frac{\mathrm{Pe}\rho}{1+\rho}e,\qquad
\bar v_{V\Theta}
=\frac{\mathrm{Pe}\rho}{(1+\rho)(1+M\rho)}e.
\]

Their raw MSD is therefore ballistic exactly when \(\mathrm{Pe}\rho>0\), even though their centered fluctuations are diffusive.



\section{Exact 21-pair witness partition}\label{s4.-exact-21-pair-witness-table}

Assume in this section that

\[
M>0,\qquad \rho>0,\qquad \mathrm{Pe}\ge0,
\]

and that parameters and the common initial state are matched. The following three cases are disjoint and give the exact partition \(12 + 6 + 3 = 21=\binom72\) unordered pairs.

\begin{enumerate}
\item \textbf{Different position-reset indicator (12 pairs):}
\[
\begin{gathered}
P/V,\ P/\Theta,\ P/V\Theta;\quad PV/V,\ PV/\Theta,\ PV/V\Theta;\\
P\Theta/V,\ P\Theta/\Theta,\ P\Theta/V\Theta;\\
PV\Theta/V,\ PV\Theta/\Theta,\ PV\Theta/V\Theta.
\end{gathered}
\]
The position-reset member has finite stationary centered variance and bounded raw MSD; the non-position member has linearly growing centered covariance. This witness also holds at \(\mathrm{Pe}=0\).

\item \textbf{Same position status, different orientation-reset indicator (6 pairs):}
\[
P/P\Theta,\quad P/PV\Theta,\quad PV/P\Theta,\quad PV/PV\Theta,\quad V/\Theta,\quad V/V\Theta.
\]
At \(\rho>0\), the orientation-reset member has \(\bar u=\rho e/(1+\rho)\) and \(U=\operatorname{diag}((\rho+2)/(\rho+4),2/(\rho+4))\); the other member has \(\bar u=0\) and \(U=I/2\). This witness is independent of \(\mathrm{Pe}\).

\item \textbf{Same position and orientation status, different velocity-reset indicator (3 pairs):}
\(P/PV, P\Theta/PV\Theta, \Theta/V\Theta\).
For \(P/PV\), each \(S_{ii}\) has the strictly positive difference shown below. For the other two pairs, \(\bar v\) differs when \(\mathrm{Pe}>0\); at \(\mathrm{Pe}=0\), each diagonal velocity moment differs by \(\rho/(2+M\rho)>0\).
\end{enumerate}

For the \(P/PV\) pair, the all-activity witness is

\[
\begin{aligned}
S_{P,ii}-S_{PV,ii}
=
\frac{\rho\left[
M^2\mathrm{Pe}^2\rho+M^2\mathrm{Pe}^2+2M^2\rho+2M^2
+3M\mathrm{Pe}^2+2M\rho+4M+2\right]}
{2(M+1)(M\rho+2)(M\rho+M+1)}
>0 .
\end{aligned}
\]

For the two orientation-reset pairs $P\Theta/PV\Theta$ and $\Theta/V\Theta$, the mean-velocity witness just cited applies only at $\mathrm{Pe}>0$. An all-activity witness is nevertheless available in the kinetic trace itself, and it is common to both pairs:
\[
\begin{aligned}
&\operatorname{Tr}S_{P\Theta}-\operatorname{Tr}S_{PV\Theta}
=\operatorname{Tr}S_{\Theta}-\operatorname{Tr}S_{V\Theta}\\
&=\frac{\rho\left[
M^2\mathrm{Pe}^2\rho^2+3M\mathrm{Pe}^2\rho+M\mathrm{Pe}^2
+2M\rho^2+4M\rho+2M+2\rho+2\right]}
{(\rho+1)(M\rho+2)(M\rho+M+1)}
>0 ,
\end{aligned}
\]
every numerator term being nonnegative on $M>0$, $\rho>0$, $\mathrm{Pe}\ge0$ with strictly positive constant contributions, and every denominator factor positive. Together with the $P/PV$ gap above, this supplies a strictly positive kinetic-trace separation for all three velocity pairs at every point of the domain, which is the form used by the three-observable decoder of the main article; the mean-velocity witness is then no longer needed for those pairs.

Because every unordered pair occurs in exactly one row, these witnesses prove

\[
\boxed{
\mathcal R_2(X;M,\mathrm{Pe},\rho)
\ne
\mathcal R_2(Y;M,\mathrm{Pe},\rho)
\quad\text{for every }X\ne Y
}
\]

on the stated finite-parameter domain. The conclusion is full-record distinguishability, not practical identifiability from finite noisy data.



\section{Restricted observational identities}\label{s5.-restricted-identity-ledger}

The following exact equalities concern a named sector, contraction, condition, resonance, or limit. None is equality of the full records. Except for the full internal-process identities in S5.1, stationary or long-time quantities are meant where a time argument is suppressed.

\subsection{Internal-sector identities}\label{s5.1-internal-sector-identities}

For every finite physical parameter point,

\[
PV\equiv V,\qquad
P\Theta\equiv\Theta,\qquad
PV\Theta\equiv V\Theta
\quad\text{on the full internal }(v,u)\text{ process}.
\]

Position reset is absent from the \((v,u)\) generator. Hence all internal means, \(S,W,U\), centered internal blocks, and higher internal statistics match for each pair. Their spatial records do not match: the left member is localized and the right member transports.

\subsection{Contracted-observable degeneracies}\label{s5.2-contracted-observable-degeneracies}

For \(PV/PV\Theta\), the named scalar triple \((\mathrm{TrS},\mathrm{TrC},\mathrm{TrR})\) agrees exactly:

\[
\boxed{
\operatorname{Tr}S_{PV}=\operatorname{Tr}S_{PV\Theta},\quad
\operatorname{Tr}C_{PV}=\operatorname{Tr}C_{PV\Theta},\quad
\operatorname{Tr}R_{PV}=\operatorname{Tr}R_{PV\Theta}.
}
\]

The equality follows from \(\operatorname{Tr}U=1\). It does not extend to centered spatial variance; in the active positive-rate interior,

\[
\operatorname{Tr}\operatorname{Cov}_{PV}(r)
-\operatorname{Tr}\operatorname{Cov}_{PV\Theta}(r)
=\frac{\mathrm{Pe}^2}
{(1+\rho)^2(1+M\rho)^2}>0.
\]

For \(V/V\Theta\),

\[
\boxed{\operatorname{Tr}S_V=\operatorname{Tr}S_{V\Theta}}
\]

for all finite physical parameters. At \(\rho>0\), orientation moments separate the records; in the active interior, \(V\Theta\) also drifts and
\[
D_{\mathrm{eff},V}-D_{\mathrm{eff},V\Theta}
=\frac{\mathrm{Pe}^2\rho\,P(M,\rho)}
{2(1+\rho)^3(1+M\rho)^3D_\star}>0,
\]
where \(D_\star=1+M(1+\rho)\) and
\[
P(M,\rho)=M^2(\rho^4+4\rho^3+4\rho^2+\rho)
+M(2\rho^3+7\rho^2+4\rho)+\rho^2+3\rho+1.
\]
All coefficients of \(P\) are positive for \(M,\rho>0\).

\subsection{Resonant identity}\label{s5.3-resonant-identity}

For \(P/P\Theta\), each member of the named triple \((\mathrm{TrS},\mathrm{TrC},\mathrm{TrR})\),

\[
\boxed{(\operatorname{Tr}S,\operatorname{Tr}C,\operatorname{Tr}R)}
\]

has a protocol difference proportional to \(\mathrm{Pe}^2(M\rho-1)\). The triple therefore agrees at the passive endpoint \(\mathrm{Pe}=0\) and on the active resonance

\[
\boxed{M\rho=1.}
\]

The full records remain distinct. For example, at \(M=\rho=1\) and \(\mathrm{Pe}=6/5\), \(\bar u_x\) is \(0\) for \(P\) and \(1/2\) for \(P\Theta\), while \(\bar r_x\) is \(0\) and \(3/5\), respectively.

\subsection{Conditional identities}\label{s5.4-conditional-identities}

At \(\mathrm{Pe}=0\), the translational sectors satisfy

\[
P=P\Theta,\qquad
PV=PV\Theta,\qquad
V=V\Theta,
\]

but orientation outputs still separate each pair for \(\rho>0\).

The orientation-only diffusion coefficients recover Eq.~(46) of Kumar, Sadekar and Basu~\cite{KumarSadekarBasu}, after setting the rotational diffusivity to unity and adding the independent translational contribution. Translational inertia leaves these long-time coefficients unchanged when velocity is retained; the result here is their comparison with velocity-reset transport.

For \(\mathrm{Pe}>0\) and \(\rho>0\), \(D_\Theta\) is isotropic only at \(\rho=1/2\); it can therefore equal the isotropic tensor \(D_V\) only there. On that line, tensor equality reduces to one scalar equation, which gives the exact curve

\[
\boxed{
\rho=\frac12,\qquad
\mathrm{Pe}^2=
\frac{27M(M+6)(3M+2)}
{2(44-9M-80M^2-12M^3)},\qquad
0<M<M_\star
}
\]

where \(M_\star\) is the unique positive root of \(44-9M-80M^2-12M^3=0\) (\(M_\star=0.658170892\ldots\)). The denominator is positive on the stated domain. This is not a full-record identity: \(\Theta\) has nonzero mean drift and ballistic raw MSD, whereas \(V\) has zero drift and diffusive raw MSD. The curve is an exact analytic identity. At the passive endpoint \(\mathrm{Pe}=0\), \(D_\Theta=I\) for every positive \(\rho\) (Sec.~S6.3), so the positive-activity isotropy statement does not apply. Three points were simulated for the effective diffusion coefficient and mean drift (Table S1).

\subsection{Limiting identities}\label{s5.5-limiting-identities}

At fixed \(\rho>0\) and \(\mathrm{Pe}\ge0\), first derive the exact finite-\(M\) stationary position moments (for position-reset protocols) or long-time drift/diffusion coefficients (for non-position protocols), and then take the position-output limit \(M\to0^+\). In this order,

\[
P=PV,\qquad
P\Theta=PV\Theta,\qquad
\Theta=V\Theta
\]

for the displayed position moments and transport coefficients.

At \(\rho=0\), the jump term vanishes and all seven protocols reduce to the same free process. The \(M\to0^+\) statement does not assert convergence of the velocity record, a transient path-space topology, the reverse order, or a joint reset-rate limit. Equal rare- or frequent-reset leading coefficients are identities only of those named coefficients, not of the finite-rate records.



\section{Physical baselines and limiting regimes}\label{s6.-mandatory-baselines-and-limits}

\subsection{Complete-reset baseline}\label{s6.1-complete-reset-baseline}

For \(PV\Theta\), the stationary means reduce to

\[
\bar u=\frac{\rho}{1+\rho}e,\qquad
\bar v=\frac{\rho\,\mathrm{Pe}}
{(1+\rho)(1+M\rho)}e,\qquad
\bar r=\frac{\mathrm{Pe}}
{(1+\rho)(1+M\rho)}e.
\]

The component solution from S3.1 reproduces the complete-reset inertial active-particle baseline of Patel and Shee~\cite{PatelShee}.

\subsection{Zero-rate and rare-reset limit}\label{s6.2-zero-rate-and-rare-reset-limit}

For every position-reset protocol,

\[
\lim_{\rho\to0^+}\operatorname{Tr}S
=\frac{2}{M}+\frac{\mathrm{Pe}^2}{1+M},
\]

\[
\lim_{\rho\to0^+}\rho\,\operatorname{Tr}R
=4+2\mathrm{Pe}^2
=4D_{\mathrm{eff,free}},
\qquad
D_{\mathrm{eff,free}}=1+\frac{\mathrm{Pe}^2}{2}.
\]

For each non-position protocol,

\[
\lim_{\rho\to0^+}\bar v=0,\qquad
\lim_{\rho\to0^+}\operatorname{Tr}S
=\frac{2}{M}+\frac{\mathrm{Pe}^2}{1+M},
\]

\[
\lim_{\rho\to0^+}D
=\left(1+\frac{\mathrm{Pe}^2}{2}\right)I.
\]

The order matters for position resetting. At each fixed \(\rho>0\), first taking \(t\to\infty\) produces a stationary position law; its variance then diverges as \(\rho\to0^+\). Setting \(\rho=0\) first removes resetting and produces free diffusion rather than a stationary law. The scaled limit above connects these two descriptions without asserting that the two limits commute.

\subsection{Passive underdamped limit}\label{s6.3-passive-underdamped-limit}

At \(\mathrm{Pe}=0\), orientation reset drops out of translational moments. For position-reset protocols,

\[
\begin{array}{c|cc}
\text{class} & S & R\\ \hline
P=P\Theta
& I/M
& \dfrac{2I}{\rho(1+M\rho)}\\[3mm]
PV=PV\Theta
& \dfrac{2I}{M(2+M\rho)}
& \dfrac{4I}{\rho(1+M\rho)(2+M\rho)}.
\end{array}
\]

For non-position protocols,

\[
D_\Theta=I,\qquad
D_V=D_{V\Theta}
=\frac{2I}{(1+M\rho)(2+M\rho)}.
\]

These are translational identities only; orientation records remain distinct when \(\rho>0\).

\subsection{Singular overdamped position-process limit}\label{s6.4-singular-overdamped-position-process-limit}

The \(M\to0^+\) comparison is a singular position-process limit. The limiting equation is

\[
dr=\mathrm{Pe}\,u\,dt+\sqrt2\,dW_t.
\]

For the localized protocols, the procedure is: solve the exact finite-\(M\) stationary position moments at fixed \((\rho,\mathrm{Pe})\), then take \(M\to0^+\). This gives

\[
P=PV,\qquad P\Theta=PV\Theta
\quad\text{in position second moments},
\]

\[
R_{\mathrm{od}}
=\frac{2I}{\rho}
+\frac{2\mathrm{Pe}^2U}{\rho(1+\rho)},
\]

\[
\bar r=
\begin{cases}
0,&P,PV,\\[1mm]
\mathrm{Pe}\,e/(1+\rho),&P\Theta,PV\Theta,
\end{cases}
\qquad
\operatorname{Cov}(r)=R_{\mathrm{od}}-\bar r\bar r^\mathsf T,
\]

where \(U=I/2\) without orientation reset and

\[
U=\operatorname{diag}\left(\frac{\rho+2}{\rho+4},
                           \frac{2}{\rho+4}\right)
\]

with orientation reset.

For the non-position protocols, compute the finite-\(M\) diffusion tensor first and then take \(M\to0^+\):

\[
D_V\longrightarrow
\left(1+\frac{\mathrm{Pe}^2}{2}\right)I,
\qquad
D_\Theta\longrightarrow
I+\frac{\mathrm{Pe}^2}{1+\rho}\Sigma_u,
\qquad
D_{V\Theta}\longrightarrow
I+\frac{\mathrm{Pe}^2}{1+\rho}\Sigma_u.
\]

Velocity is not an independent state at \(M=0\). These statements therefore assert convergence of the displayed position outputs, not convergence of the finite-inertia velocity record. No direct substitution \(M=0\) is made in a velocity formula. Any joint limit involving \(M\), \(\rho\), or \(t\to\infty\) must state its order; the results above do not assert that those limits commute.

\subsection{Frequent-reset asymptotics}\label{s6.5-frequent-reset-asymptotics}

For position-reset protocols,

\[
\begin{aligned}
\lim_{\rho\to\infty}\rho^2E_P|r|^2
&=\frac4M+\frac{2\mathrm{Pe}^2}{1+M},\\
\lim_{\rho\to\infty}\rho^2E_{P\Theta}|r|^2
&=\frac4M+2\mathrm{Pe}^2,\\
\lim_{\rho\to\infty}\rho^3E_{PV}|r|^2
=\lim_{\rho\to\infty}\rho^3E_{PV\Theta}|r|^2
&=\frac8{M^2}.
\end{aligned}
\]

For \(P\Theta\), \(\lim_{\rho\to\infty}\rho^2|\bar r|^2=\mathrm{Pe}^2\), so its centered variance coefficient is \(4/M+\mathrm{Pe}^2\).

For non-position protocols,

\[
\begin{aligned}
\lim_{\rho\to\infty}\rho^2D_{\mathrm{eff},V}
&=\frac{\mathrm{Pe}^2+4}{2M^2},\\
\lim_{\rho\to\infty}D_{\mathrm{eff},\Theta}
&=1,\qquad
\lim_{\rho\to\infty}\bar v_\Theta=\mathrm{Pe}\,e,\\
\lim_{\rho\to\infty}\rho^2D_{\mathrm{eff},V\Theta}
&=\frac2{M^2},\qquad
\lim_{\rho\to\infty}\rho\bar v_{V\Theta}
=\frac{\mathrm{Pe}}{M}e.
\end{aligned}
\]

For every fixed finite \(\rho>0\), \(\Theta\) and \(V\Theta\) have a ballistic raw-MSD term whenever \(\mathrm{Pe}>0\). The vanishing \(V\Theta\) drift coefficient as \(\rho\to\infty\) does not retroactively change that finite-rate classification. This is another noncommuting classification/asymptotic distinction: a leading coefficient may vanish in a parameter limit although the finite-parameter growth exponent remains ballistic.

\section{Scope of the exact results}\label{s7.-claim-boundary}

The exact result is full-record second-order distinguishability of all seven nonempty reset protocols for \(M>0\), \(\rho>0\), and \(\mathrm{Pe}\ge0\), together with a scope-labeled account of information loss under internal-sector restriction, scalar contraction, parameter-specific conditions, resonance, and singular limits. Trace equality is not matrix equality; raw-MSD equality is not centered-covariance equality; and none of the restricted identities is full protocol equivalence.

The supplement makes no practical-identifiability claim and introduces no higher-order or distributional claim, apart from the exact internal-process invariance in Sec.~S5.1.
